\documentclass[aspectratio=169, 11pt]{beamer}
% ============================================================================
% Preambulo compartido para presentaciones Beamer
% Estadistica 2026-II - Pontificia Universidad Javeriana
% Incluir con: % ============================================================================
% Preambulo compartido para presentaciones Beamer
% Estadistica 2026-II - Pontificia Universidad Javeriana
% Incluir con: % ============================================================================
% Preambulo compartido para presentaciones Beamer
% Estadistica 2026-II - Pontificia Universidad Javeriana
% Incluir con: \input{../preambulo_beamer.tex} despues de \documentclass
% ============================================================================

\usepackage[T1]{fontenc}
\usepackage[utf8]{inputenc}
\usepackage[spanish]{babel}
\usepackage{amsmath, amssymb, amsfonts}
\usepackage{booktabs}
\usepackage{graphicx}
\usepackage{tikz}
\usetikzlibrary{shapes.geometric, positioning, arrows.meta, calc}
\usepackage{pgfplots}
\pgfplotsset{compat=1.18}
\usepackage{listings}
\usepackage{xcolor}
\usepackage{hyperref}
\usepackage{multicol}

% --- Tema Beamer (minimalista) ---
\usetheme{default}
\usecolortheme{default}
\usefonttheme{default}
\setbeamertemplate{navigation symbols}{}
\setbeamertemplate{caption}[numbered]
\setbeamertemplate{footline}{%
  \hspace{1em}{\tiny\url{https://github.com/polanco-jaime/Curso-Estadistica-2026-3}}\hfill\usebeamerfont{footline}\insertframenumber/\inserttotalframenumber\hspace{1em}\vspace{0.5em}%
}
\setbeamertemplate{itemize items}[circle]
\setbeamertemplate{enumerate items}[default]

% --- Colores institucionales (sutiles) ---
\definecolor{javeriana}{RGB}{0, 62, 105}
\definecolor{javerianalight}{RGB}{0, 105, 170}
\definecolor{codigobg}{RGB}{245, 245, 245}
\definecolor{codigofg}{RGB}{40, 40, 40}
\definecolor{comentario}{RGB}{100, 100, 100}
\definecolor{keyword}{RGB}{0, 100, 150}
\definecolor{string}{RGB}{180, 60, 30}

\setbeamercolor{title}{fg=javeriana}
\setbeamercolor{frametitle}{fg=javeriana}
\setbeamercolor{structure}{fg=javeriana}
\setbeamercolor{block title}{fg=white, bg=javeriana!75}
\setbeamercolor{block body}{bg=javeriana!5}
\setbeamercolor{block title alerted}{fg=white, bg=red!70!black}
\setbeamercolor{block body alerted}{bg=red!5}
\setbeamercolor{block title example}{fg=white, bg=green!40!black}
\setbeamercolor{block body example}{bg=green!5}
\setbeamercolor{item}{fg=javeriana}

\setbeamerfont{title}{series=\bfseries, size=\Large}
\setbeamerfont{frametitle}{series=\bfseries}

% --- Configuracion de listings para R ---
\lstset{
  language=R,
  basicstyle=\ttfamily\footnotesize,
  backgroundcolor=\color{codigobg},
  commentstyle=\color{comentario}\itshape,
  keywordstyle=\color{keyword}\bfseries,
  stringstyle=\color{string},
  numbers=none,
  frame=single,
  rulecolor=\color{javeriana!30},
  breaklines=true,
  showstringspaces=false,
  columns=flexible,
  morekeywords={library, ggplot, aes, geom_histogram, geom_boxplot,
    geom_point, geom_smooth, geom_density, geom_bar, geom_line,
    geom_vline, geom_hline, geom_errorbar, labs, theme_minimal,
    facet_wrap, scale_fill_manual, scale_color_manual,
    summarise, group_by, filter, mutate, select, arrange,
    read.csv, head, str, summary, mean, median, sd, var, cor,
    lm, t.test, chisq.test, wilcox.test, kruskal.test, aov,
    pnorm, qnorm, dnorm, rnorm, qt, pt, qchisq, pchisq, qf, pf,
    confint, coeftest, vcovHC, bptest, vif, cooks.distance,
    TukeyHSD, table, prop.table}
}

% --- Comandos personalizados ---
\newcommand{\R}{\texttt{R}}
\newcommand{\code}[1]{\texttt{\footnotesize #1}}
\newcommand{\variable}[1]{\texttt{#1}}
\newcommand{\paquete}[1]{\texttt{#1}}

% --- Informacion del curso ---
\institute{Pontificia Universidad Javeriana\\
  Facultad de Ciencias Econ\'omicas y Administrativas}
\date{2026-II}
 despues de \documentclass
% ============================================================================

\usepackage[T1]{fontenc}
\usepackage[utf8]{inputenc}
\usepackage[spanish]{babel}
\usepackage{amsmath, amssymb, amsfonts}
\usepackage{booktabs}
\usepackage{graphicx}
\usepackage{tikz}
\usetikzlibrary{shapes.geometric, positioning, arrows.meta, calc}
\usepackage{pgfplots}
\pgfplotsset{compat=1.18}
\usepackage{listings}
\usepackage{xcolor}
\usepackage{hyperref}
\usepackage{multicol}

% --- Tema Beamer (minimalista) ---
\usetheme{default}
\usecolortheme{default}
\usefonttheme{default}
\setbeamertemplate{navigation symbols}{}
\setbeamertemplate{caption}[numbered]
\setbeamertemplate{footline}{%
  \hspace{1em}{\tiny\url{https://github.com/polanco-jaime/Curso-Estadistica-2026-3}}\hfill\usebeamerfont{footline}\insertframenumber/\inserttotalframenumber\hspace{1em}\vspace{0.5em}%
}
\setbeamertemplate{itemize items}[circle]
\setbeamertemplate{enumerate items}[default]

% --- Colores institucionales (sutiles) ---
\definecolor{javeriana}{RGB}{0, 62, 105}
\definecolor{javerianalight}{RGB}{0, 105, 170}
\definecolor{codigobg}{RGB}{245, 245, 245}
\definecolor{codigofg}{RGB}{40, 40, 40}
\definecolor{comentario}{RGB}{100, 100, 100}
\definecolor{keyword}{RGB}{0, 100, 150}
\definecolor{string}{RGB}{180, 60, 30}

\setbeamercolor{title}{fg=javeriana}
\setbeamercolor{frametitle}{fg=javeriana}
\setbeamercolor{structure}{fg=javeriana}
\setbeamercolor{block title}{fg=white, bg=javeriana!75}
\setbeamercolor{block body}{bg=javeriana!5}
\setbeamercolor{block title alerted}{fg=white, bg=red!70!black}
\setbeamercolor{block body alerted}{bg=red!5}
\setbeamercolor{block title example}{fg=white, bg=green!40!black}
\setbeamercolor{block body example}{bg=green!5}
\setbeamercolor{item}{fg=javeriana}

\setbeamerfont{title}{series=\bfseries, size=\Large}
\setbeamerfont{frametitle}{series=\bfseries}

% --- Configuracion de listings para R ---
\lstset{
  language=R,
  basicstyle=\ttfamily\footnotesize,
  backgroundcolor=\color{codigobg},
  commentstyle=\color{comentario}\itshape,
  keywordstyle=\color{keyword}\bfseries,
  stringstyle=\color{string},
  numbers=none,
  frame=single,
  rulecolor=\color{javeriana!30},
  breaklines=true,
  showstringspaces=false,
  columns=flexible,
  morekeywords={library, ggplot, aes, geom_histogram, geom_boxplot,
    geom_point, geom_smooth, geom_density, geom_bar, geom_line,
    geom_vline, geom_hline, geom_errorbar, labs, theme_minimal,
    facet_wrap, scale_fill_manual, scale_color_manual,
    summarise, group_by, filter, mutate, select, arrange,
    read.csv, head, str, summary, mean, median, sd, var, cor,
    lm, t.test, chisq.test, wilcox.test, kruskal.test, aov,
    pnorm, qnorm, dnorm, rnorm, qt, pt, qchisq, pchisq, qf, pf,
    confint, coeftest, vcovHC, bptest, vif, cooks.distance,
    TukeyHSD, table, prop.table}
}

% --- Comandos personalizados ---
\newcommand{\R}{\texttt{R}}
\newcommand{\code}[1]{\texttt{\footnotesize #1}}
\newcommand{\variable}[1]{\texttt{#1}}
\newcommand{\paquete}[1]{\texttt{#1}}

% --- Informacion del curso ---
\institute{Pontificia Universidad Javeriana\\
  Facultad de Ciencias Econ\'omicas y Administrativas}
\date{2026-II}
 despues de \documentclass
% ============================================================================

\usepackage[T1]{fontenc}
\usepackage[utf8]{inputenc}
\usepackage[spanish]{babel}
\usepackage{amsmath, amssymb, amsfonts}
\usepackage{booktabs}
\usepackage{graphicx}
\usepackage{tikz}
\usetikzlibrary{shapes.geometric, positioning, arrows.meta, calc}
\usepackage{pgfplots}
\pgfplotsset{compat=1.18}
\usepackage{listings}
\usepackage{xcolor}
\usepackage{hyperref}
\usepackage{multicol}

% --- Tema Beamer (minimalista) ---
\usetheme{default}
\usecolortheme{default}
\usefonttheme{default}
\setbeamertemplate{navigation symbols}{}
\setbeamertemplate{caption}[numbered]
\setbeamertemplate{footline}{%
  \hspace{1em}{\tiny\url{https://github.com/polanco-jaime/Curso-Estadistica-2026-3}}\hfill\usebeamerfont{footline}\insertframenumber/\inserttotalframenumber\hspace{1em}\vspace{0.5em}%
}
\setbeamertemplate{itemize items}[circle]
\setbeamertemplate{enumerate items}[default]

% --- Colores institucionales (sutiles) ---
\definecolor{javeriana}{RGB}{0, 62, 105}
\definecolor{javerianalight}{RGB}{0, 105, 170}
\definecolor{codigobg}{RGB}{245, 245, 245}
\definecolor{codigofg}{RGB}{40, 40, 40}
\definecolor{comentario}{RGB}{100, 100, 100}
\definecolor{keyword}{RGB}{0, 100, 150}
\definecolor{string}{RGB}{180, 60, 30}

\setbeamercolor{title}{fg=javeriana}
\setbeamercolor{frametitle}{fg=javeriana}
\setbeamercolor{structure}{fg=javeriana}
\setbeamercolor{block title}{fg=white, bg=javeriana!75}
\setbeamercolor{block body}{bg=javeriana!5}
\setbeamercolor{block title alerted}{fg=white, bg=red!70!black}
\setbeamercolor{block body alerted}{bg=red!5}
\setbeamercolor{block title example}{fg=white, bg=green!40!black}
\setbeamercolor{block body example}{bg=green!5}
\setbeamercolor{item}{fg=javeriana}

\setbeamerfont{title}{series=\bfseries, size=\Large}
\setbeamerfont{frametitle}{series=\bfseries}

% --- Configuracion de listings para R ---
\lstset{
  language=R,
  basicstyle=\ttfamily\footnotesize,
  backgroundcolor=\color{codigobg},
  commentstyle=\color{comentario}\itshape,
  keywordstyle=\color{keyword}\bfseries,
  stringstyle=\color{string},
  numbers=none,
  frame=single,
  rulecolor=\color{javeriana!30},
  breaklines=true,
  showstringspaces=false,
  columns=flexible,
  morekeywords={library, ggplot, aes, geom_histogram, geom_boxplot,
    geom_point, geom_smooth, geom_density, geom_bar, geom_line,
    geom_vline, geom_hline, geom_errorbar, labs, theme_minimal,
    facet_wrap, scale_fill_manual, scale_color_manual,
    summarise, group_by, filter, mutate, select, arrange,
    read.csv, head, str, summary, mean, median, sd, var, cor,
    lm, t.test, chisq.test, wilcox.test, kruskal.test, aov,
    pnorm, qnorm, dnorm, rnorm, qt, pt, qchisq, pchisq, qf, pf,
    confint, coeftest, vcovHC, bptest, vif, cooks.distance,
    TukeyHSD, table, prop.table}
}

% --- Comandos personalizados ---
\newcommand{\R}{\texttt{R}}
\newcommand{\code}[1]{\texttt{\footnotesize #1}}
\newcommand{\variable}[1]{\texttt{#1}}
\newcommand{\paquete}[1]{\texttt{#1}}

% --- Informacion del curso ---
\institute{Pontificia Universidad Javeriana\\
  Facultad de Ciencias Econ\'omicas y Administrativas}
\date{2026-II}

\title{Sesión 10: Regresión Lineal Múltiple}
\subtitle{Múltiples regresores, OVB y multicolinealidad}
\author{Prof.\ Jaime Polanco Jim\'enez}
\date{Vie 30 oct 2026}

\begin{document}

\begin{frame}
\titlepage
\end{frame}

\begin{frame}{Agenda}
\tableofcontents
\end{frame}

% ============================================================================
\section{Modelo de regresión múltiple}
% ============================================================================

\begin{frame}{M\'as all\'a de un solo predictor}
\begin{block}{Limitaci\'on de la regresi\'on simple}
En regresi\'on simple ($Y = \beta_0 + \beta_1 X + \varepsilon$), solo consideramos \textbf{un} predictor. Pero en la pr\'actica, m\'ultiples factores influyen en $Y$ simult\'aneamente.
\end{block}

\vspace{0.2cm}
\textbf{Ejemplo:} El puntaje de ingl\'es en Saber Pro puede depender de:
\begin{itemize}
  \item Puntaje en Saber 11, estrato socioecon\'omico, tipo de IES
  \item Nivel educativo de los padres, regi\'on del pa\'is, ingresos del hogar
\end{itemize}

\vspace{0.2cm}
\begin{exampleblock}{Objetivo de regresi\'on m\'ultiple}
Estimar el efecto de cada predictor \textbf{controlando por} (manteniendo constantes) los dem\'as predictores.

\smallskip
\textbf{Analog\'ia:} Quieren predecir su nota del parcial controlando por horas de estudio, screen time diario, y si trabajan. La regresi\'on m\'ultiple permite aislar el efecto de cada factor.
\end{exampleblock}
\end{frame}

\begin{frame}{El modelo de regresión lineal múltiple}
\begin{block}{Modelo poblacional}
\[
Y_i = \beta_0 + \beta_1 X_{1i} + \beta_2 X_{2i} + \cdots + \beta_p X_{pi} + \varepsilon_i
\]
donde:
\begin{itemize}
  \item $Y_i$ = variable dependiente (respuesta)
  \item $X_{1i}, X_{2i}, \ldots, X_{pi}$ = variables independientes (predictores, regresores)
  \item $\beta_0$ = intercepto
  \item $\beta_j$ = coeficiente de $X_j$ ($j = 1, 2, \ldots, p$)
  \item $\varepsilon_i$ = error aleatorio
\end{itemize}
\end{block}

\vspace{0.2cm}
\textbf{Interpretación de $\beta_j$:}\\
$\beta_j$ representa el cambio esperado en $Y$ cuando $X_j$ aumenta en una unidad, \textbf{manteniendo constantes todas las demás variables} (ceteris paribus).
\end{frame}

\begin{frame}{Ecuación estimada}
Estimamos los parámetros $\beta_0, \beta_1, \ldots, \beta_p$ usando OLS (mínimos cuadrados ordinarios), obteniendo:
\[
\hat{Y}_i = \hat{\beta}_0 + \hat{\beta}_1 X_{1i} + \hat{\beta}_2 X_{2i} + \cdots + \hat{\beta}_p X_{pi}
\]

\vspace{0.2cm}
\textbf{Residuos:}
\[
e_i = Y_i - \hat{Y}_i
\]

\vspace{0.2cm}
\textbf{OLS minimiza:}
\[
\text{SSE} = \sum_{i=1}^{n} e_i^2 = \sum_{i=1}^{n} (Y_i - \hat{Y}_i)^2
\]

\vspace{0.2cm}
\begin{alertblock}{Interpretación matricial}
En álgebra matricial: $\boldsymbol{\hat{\beta}} = (\mathbf{X}'\mathbf{X})^{-1}\mathbf{X}'\mathbf{Y}$\\
(No entraremos en detalle, pero R usa esta fórmula internamente)
\end{alertblock}
\end{frame}

\begin{frame}[fragile]{Sintaxis en R: \texttt{lm(Y \textasciitilde{} X1 + X2 + ...)}}
\begin{lstlisting}[basicstyle=\ttfamily\scriptsize]
# Cargar datos ICFES para Negocios Internacionales
icfes_ni <- read.csv("datos/icfes_ni_2025.csv")

# Regresion multiple: MOD_INGLES_PUNT como funcion de
# varias variables
modelo_mult <- lm(MOD_INGLES_PUNT ~
                    INGLES_SABER11_PUNT +
                    ESTRATO +
                    PRIVADA +
                    EDUCACION_PADRE,
                  data = icfes_ni)

# Ver coeficientes
coef(modelo_mult)

#        (Intercept) INGLES_SABER11_PUNT             ESTRATO
#              32.45                1.67                2.13
#            PRIVADA     EDUCACION_PADRE
#               3.82                1.25
\end{lstlisting}
\end{frame}

\begin{frame}{Interpretación de los coeficientes}
Del modelo anterior:
\[
\widehat{\text{MOD\_INGLES}} = 32.45 + 1.67 \times \text{Saber11} + 2.13 \times \text{Estrato} + 3.82 \times \text{Privada} + 1.25 \times \text{EducPadre}
\]

\vspace{0.2cm}
\textbf{Interpretaciones ceteris paribus:}
\begin{itemize}
  \item $\hat{\beta}_1 = 1.67$: Por cada punto adicional en Saber 11, el puntaje de inglés en Saber Pro aumenta \textbf{en promedio} 1.67 puntos, \textbf{manteniendo constantes} estrato, tipo de IES y educación del padre.

  \item $\hat{\beta}_2 = 2.13$: Por cada nivel adicional de estrato, el puntaje aumenta 2.13 puntos, \textbf{controlando por} los demás factores.

  \item $\hat{\beta}_3 = 3.82$: Los estudiantes de IES privadas obtienen 3.82 puntos más que los de públicas, \textbf{en promedio}, controlando por las demás variables.

  \item $\hat{\beta}_4 = 1.25$: Por cada año adicional de educación del padre, el puntaje aumenta 1.25 puntos (ceteris paribus).
\end{itemize}
\end{frame}

% ============================================================================
\section{Ejemplo: modelo completo del inglés en NI}
% ============================================================================

\begin{frame}[fragile]{Modelo extendido con transformaciones}
\begin{lstlisting}[basicstyle=\ttfamily\scriptsize]
# Crear variables dummy para region (referencia: Andina)
icfes_ni$REGION_CARIBE <- ifelse(icfes_ni$REGION == "Caribe", 1, 0)
icfes_ni$REGION_PACIFICA <- ifelse(icfes_ni$REGION == "Pacifica", 1, 0)
icfes_ni$REGION_ORINOQUIA <- ifelse(icfes_ni$REGION == "Orinoquia", 1, 0)
icfes_ni$REGION_BOGOTA <- ifelse(icfes_ni$REGION == "Bogota", 1, 0)

# Transformacion logaritmica de ingresos (reduce asimetria)
icfes_ni$LOG_INGRESOS <- log(icfes_ni$INGRESOS + 1)

# Modelo completo
modelo_completo <- lm(MOD_INGLES_PUNT ~
                        INGLES_SABER11_PUNT +
                        ESTRATO +
                        PRIVADA +
                        LOG_INGRESOS +
                        REGION_CARIBE +
                        REGION_PACIFICA +
                        REGION_ORINOQUIA +
                        REGION_BOGOTA,
                      data = icfes_ni)
\end{lstlisting}
\end{frame}

\begin{frame}[fragile]{Resultados del modelo completo}
\begin{lstlisting}
summary(modelo_completo)

# Coefficients:
#                      Estimate Std. Error t value Pr(>|t|)
# (Intercept)           28.342      2.145   13.21  < 2e-16 ***
# INGLES_SABER11_PUNT    1.612      0.025   64.48  < 2e-16 ***
# ESTRATO                1.876      0.312    6.01  2.1e-09 ***
# PRIVADA                3.254      0.672    4.84  1.3e-06 ***
# LOG_INGRESOS           2.145      0.421    5.09  3.6e-07 ***
# REGION_CARIBE         -2.543      0.834   -3.05  0.00230 **
# REGION_PACIFICA       -1.267      0.752   -1.68  0.09234 .
# REGION_ORINOQUIA      -4.123      1.234   -3.34  0.00084 ***
# REGION_BOGOTA          5.234      0.812    6.44  1.3e-10 ***
# ---
# Residual standard error: 10.85 on 5076 DF
# Multiple R-squared:  0.7145, Adjusted R-squared:  0.7140
# F-statistic: 1586 on 8 and 5076 DF,  p-value: < 2.2e-16
\end{lstlisting}
\end{frame}

\begin{frame}{Interpretación del modelo completo}
\textbf{Variables continuas:}
\begin{itemize}
  \item INGLES\_SABER11\_PUNT: $\hat{\beta} = 1.61$ (altamente significativo)\\
    Por cada punto adicional en Saber 11, +1.61 en Saber Pro (controlando por todo lo demás)

  \item ESTRATO: $\hat{\beta} = 1.88$ (significativo)\\
    Por cada nivel de estrato, +1.88 puntos

  \item LOG\_INGRESOS: $\hat{\beta} = 2.15$ (significativo)\\
    Interpretación: un aumento del 1\% en ingresos $\Rightarrow$ +0.0215 puntos\\
    (semi-elasticidad: $\beta/100$)
\end{itemize}

\vspace{0.2cm}
\textbf{Variable dummy:}
\begin{itemize}
  \item PRIVADA: $\hat{\beta} = 3.25$ (significativo)\\
    IES privadas obtienen 3.25 puntos más que públicas, en promedio, controlando por estrato, ingresos, región, etc.
\end{itemize}
\end{frame}

\begin{frame}{Interpretaci\'on de dummies regionales}
\textbf{Regi\'on de referencia:} Andina (omitida para evitar multicolinealidad perfecta)

\vspace{0.1cm}
\textbf{Coeficientes regionales (relativo a Andina):}
\begin{itemize}
  \item REGION\_CARIBE: $\hat{\beta} = -2.54$ (significativo)\\
    Estudiantes del Caribe obtienen 2.54 puntos \textbf{menos} que los de la Andina (controlando por todo)

  \item REGION\_PACIFICA: $\hat{\beta} = -1.27$ (no significativo, p = 0.092)\\
    No hay diferencia estadísticamente significativa con Andina

  \item REGION\_ORINOQUIA: $\hat{\beta} = -4.12$ (significativo)\\
    Orinoquia obtiene 4.12 puntos menos que Andina

  \item REGION\_BOGOTA: $\hat{\beta} = +5.23$ (altamente significativo)\\
    Bogotá obtiene 5.23 puntos \textbf{más} que Andina
\end{itemize}

\vspace{0.2cm}
\begin{alertblock}{Nota}
Siempre se omite una categoría de referencia en variables categóricas para evitar colinealidad perfecta (dummy variable trap).
\end{alertblock}
\end{frame}

% ============================================================================
\section{$R^2$ y $R^2$ ajustado}
% ============================================================================

\begin{frame}{$R^2$ en regresi\'on m\'ultiple}
El $R^2$ se define igual que en regresi\'on simple:
\[
R^2 = \frac{\text{SSR}}{\text{SST}} = 1 - \frac{\text{SSE}}{\text{SST}}
\]

\vspace{0.2cm}
\textbf{Problema:} $R^2$ \textbf{siempre aumenta} (o se mantiene) al agregar variables, incluso si son irrelevantes.

\vspace{0.2cm}
\begin{exampleblock}{Ejemplo}
\begin{itemize}
  \item Modelo 1: $Y \sim X_1$ $\Rightarrow$ $R^2 = 0.65$
  \item Modelo 2: $Y \sim X_1 + X_2$ $\Rightarrow$ $R^2 = 0.66$
  \item Modelo 3: $Y \sim X_1 + X_2 + X_3$ $\Rightarrow$ $R^2 = 0.662$
\end{itemize}
Incluso si $X_3$ es ruido puro, $R^2$ no disminuir\'a.
\end{exampleblock}
\end{frame}

\begin{frame}{$R^2$ en regresi\'on m\'ultiple (cont.)}
\begin{alertblock}{Consecuencia}
No podemos comparar modelos con diferente n\'umero de variables usando solo $R^2$. Necesitamos una medida que \textbf{penalice} la complejidad. Es como agregar filtros a una foto: siempre se ve ``mejor'', pero no necesariamente m\'as real.
\end{alertblock}

\vspace{0.3cm}
\textbf{Soluci\'on:} El $R^2$ ajustado ($R^2_{\text{adj}}$) penaliza por el n\'umero de variables incluidas en el modelo.
\end{frame}

\begin{frame}{$R^2$ ajustado}
\begin{block}{Definici\'on}
\[
R^2_{\text{adj}} = 1 - \frac{\text{SSE}/(n - p - 1)}{\text{SST}/(n - 1)} = 1 - \frac{\text{MSE}}{\text{MST}}
\]
donde $n$ = observaciones, $p$ = predictores, MSE = cuadrado medio del error, MST = cuadrado medio total.
\end{block}

\vspace{0.1cm}
\textbf{Propiedades:}
\begin{itemize}
  \item $R^2_{\text{adj}} \leq R^2$
  \item $R^2_{\text{adj}}$ puede disminuir si agregamos una variable irrelevante
  \item Penaliza por el n\'umero de predictores
  \item \'Util para comparar modelos con diferente n\'umero de variables
\end{itemize}
\end{frame}

\begin{frame}[fragile]{Comparaci\'on: $R^2$ vs $R^2$ ajustado}
\begin{lstlisting}[basicstyle=\ttfamily\scriptsize]
# Modelo simple
modelo1 <- lm(MOD_INGLES_PUNT ~ INGLES_SABER11_PUNT,
              data = icfes_ni)
summary(modelo1)$r.squared      # 0.6524
summary(modelo1)$adj.r.squared  # 0.6523

# Modelo multiple (4 variables)
modelo2 <- lm(MOD_INGLES_PUNT ~ INGLES_SABER11_PUNT +
                ESTRATO + PRIVADA + LOG_INGRESOS,
              data = icfes_ni)
summary(modelo2)$r.squared      # 0.6982
summary(modelo2)$adj.r.squared  # 0.6979

# Modelo completo (8 variables)
summary(modelo_completo)$r.squared      # 0.7145
summary(modelo_completo)$adj.r.squared  # 0.7140
\end{lstlisting}

\vspace{0.1cm}
{\small \textbf{Observaci\'on:} $R^2_{\text{adj}}$ es ligeramente menor que $R^2$. La diferencia crece cuando agregamos muchas variables poco \'utiles.}
\end{frame}

\begin{frame}{¿Cuándo usar $R^2$ ajustado?}
\begin{block}{Regla práctica}
\begin{itemize}
  \item Para \textbf{reportar} el ajuste de un modelo: usar ambos ($R^2$ y $R^2_{\text{adj}}$)
  \item Para \textbf{comparar} modelos con diferente número de variables: usar $R^2_{\text{adj}}$
  \item Para \textbf{seleccionar} el mejor modelo: complementar con AIC/BIC (siguiente sección)
\end{itemize}
\end{block}

\vspace{0.2cm}
\begin{alertblock}{Precaución}
$R^2_{\text{adj}}$ NO es una probabilidad ni una proporción estricta. Puede ser negativo si el modelo es muy malo (aunque en la práctica esto es raro).
\end{alertblock}

\vspace{0.2cm}
\begin{exampleblock}{Interpretación aproximada}
Si $R^2_{\text{adj}} = 0.714$, aproximadamente el 71\% de la variabilidad de $Y$ se explica por los predictores, ajustando por el número de variables.
\end{exampleblock}
\end{frame}

% ============================================================================
\section{Pruebas de significancia}
% ============================================================================

\begin{frame}{Prueba F global}
\begin{block}{Hipótesis}
\begin{itemize}
  \item $H_0$: $\beta_1 = \beta_2 = \cdots = \beta_p = 0$ (todos los coeficientes son cero; el modelo no es útil)
  \item $H_1$: Al menos un $\beta_j \neq 0$ (al menos un predictor es útil)
\end{itemize}
\end{block}

\vspace{0.2cm}
\textbf{Estadístico:}
\[
F = \frac{\text{MSR}}{\text{MSE}} = \frac{\text{SSR}/p}{\text{SSE}/(n - p - 1)} \sim F_{p, n-p-1} \quad \text{bajo } H_0
\]

\vspace{0.2cm}
\textbf{Regla de decisión:}
\begin{itemize}
  \item Rechazamos $H_0$ si $F > F_{\alpha, p, n-p-1}$ o si p-valor $< \alpha$
\end{itemize}

\vspace{0.2cm}
\begin{exampleblock}{Interpretación}
Si rechazamos $H_0$, concluimos que \textbf{al menos una} variable predictora tiene efecto significativo sobre $Y$. El modelo en conjunto es útil.
\end{exampleblock}
\end{frame}

\begin{frame}{Relación entre $F$ y $R^2$}
El estadístico $F$ se puede expresar en función de $R^2$:
\[
F = \frac{R^2 / p}{(1 - R^2)/(n - p - 1)}
\]

\vspace{0.2cm}
\textbf{Interpretación:}
\begin{itemize}
  \item $R^2$ alto $\Rightarrow$ $F$ alto $\Rightarrow$ modelo significativo
  \item $R^2$ cercano a 0 $\Rightarrow$ $F$ cercano a 0 $\Rightarrow$ modelo no útil
\end{itemize}

\vspace{0.3cm}
\begin{exampleblock}{Ejemplo}
Modelo con $R^2 = 0.714$, $p = 8$, $n = 5085$:
\[
F = \frac{0.714 / 8}{(1 - 0.714) / (5085 - 8 - 1)} = \frac{0.08925}{0.0000564} \approx 1583
\]
Con p-valor $< 2.2 \times 10^{-16}$, rechazamos $H_0$: el modelo es altamente significativo.
\end{exampleblock}
\end{frame}

\begin{frame}{Prueba t individual}
\begin{block}{Hip\'otesis para cada coeficiente}
\begin{itemize}
  \item $H_0$: $\beta_j = 0$ (el predictor $X_j$ NO tiene efecto, controlando por los dem\'as)
  \item $H_1$: $\beta_j \neq 0$ (el predictor $X_j$ s\'i tiene efecto)
\end{itemize}
\end{block}

\vspace{0.1cm}
\textbf{Estad\'istico:}
$t_j = \hat{\beta}_j / \text{SE}(\hat{\beta}_j) \sim t_{n - p - 1}$ bajo $H_0$.
Rechazamos $H_0$ si $|t_j| > t_{\alpha/2, n-p-1}$ o si p-valor $< \alpha$.

\vspace{0.1cm}
\begin{alertblock}{Diferencia con regresi\'on simple}
\small En regresi\'on m\'ultiple, la prueba t para $\beta_j$ eval\'ua el efecto de $X_j$ \textbf{controlando por todas las dem\'as variables}. El resultado puede cambiar dr\'asticamente si agregamos o quitamos predictores.
\end{alertblock}
\end{frame}

\begin{frame}[fragile]{Interpretando el output de \texttt{summary()}}
\begin{lstlisting}
summary(modelo_completo)

# Coefficients:
#                      Estimate Std. Error t value Pr(>|t|)
# (Intercept)           28.342      2.145   13.21  < 2e-16 ***
# INGLES_SABER11_PUNT    1.612      0.025   64.48  < 2e-16 ***
# ESTRATO                1.876      0.312    6.01  2.1e-09 ***
# PRIVADA                3.254      0.672    4.84  1.3e-06 ***
# LOG_INGRESOS           2.145      0.421    5.09  3.6e-07 ***
# REGION_CARIBE         -2.543      0.834   -3.05  0.00230 **
# REGION_PACIFICA       -1.267      0.752   -1.68  0.09234 .
# REGION_ORINOQUIA      -4.123      1.234   -3.34  0.00084 ***
# REGION_BOGOTA          5.234      0.812    6.44  1.3e-10 ***
\end{lstlisting}

\vspace{0.2cm}
\textbf{Conclusión:}
\begin{itemize}
  \item Todas las variables son significativas \textbf{excepto} REGION\_PACIFICA (p = 0.092)
  \item La región Pacífica no difiere significativamente de la Andina (al 5\%)
\end{itemize}
\end{frame}

\begin{frame}{Significancia conjunta vs individual}
\begin{alertblock}{Posible discrepancia}
Es posible que:
\begin{itemize}
  \item La prueba F global sea significativa, pero ning\'un coeficiente individual lo sea
  \item O viceversa (raro en la pr\'actica)
\end{itemize}
\end{alertblock}

\vspace{0.1cm}
\textbf{Causas:} multicolinealidad, tama\~no muestral peque\~no, efectos d\'ebiles pero m\'ultiples.

\vspace{0.1cm}
\begin{block}{Recomendaci\'on}
Siempre reportar ambas pruebas (F global e individuales). Si hay discrepancia, investigar multicolinealidad (VIF) y considerar eliminar variables redundantes.
\end{block}
\end{frame}

% ============================================================================
\section{Sesgo por variable omitida (OVB)}
% ============================================================================

\begin{frame}{\textquestiondown Qu\'e es el sesgo por variable omitida?}
\begin{block}{Definici\'on}
El \textbf{sesgo por variable omitida} (Omitted Variable Bias, OVB) ocurre cuando:
\begin{enumerate}
  \item Omitimos una variable relevante del modelo
  \item Esa variable est\'a correlacionada con al menos uno de los predictores incluidos
\end{enumerate}
En este caso, los coeficientes estimados est\'an \textbf{sesgados}.
\end{block}

\vspace{0.1cm}
\begin{exampleblock}{Ejemplo}
\small Si omitimos ESTRATO del modelo $\text{Ingl\'es} = \beta_0 + \beta_1 \times \text{Privada} + \varepsilon$, y ESTRATO est\'a correlacionado con PRIVADA, entonces $\hat{\beta}_1$ \textbf{sobre-estima} el efecto de IES privada.

\smallskip
\textbf{Analog\'ia:} Si no controlan por nivel socioecon\'omico, podr\'ian pensar que tener iPad mejora las notas, cuando en realidad es que familias de mayores ingresos compran iPads \textbf{y adem\'as} invierten m\'as en educaci\'on.
\end{exampleblock}
\end{frame}

\begin{frame}{Fórmula del sesgo}
Si el modelo verdadero es:
\[
Y = \beta_0 + \beta_1 X_1 + \beta_2 X_2 + \varepsilon
\]

Pero estimamos el modelo omitido:
\[
Y = \gamma_0 + \gamma_1 X_1 + u
\]

Entonces el sesgo de $\hat{\gamma}_1$ es:
\[
E(\hat{\gamma}_1) = \beta_1 + \beta_2 \times \frac{\text{Cov}(X_1, X_2)}{\text{Var}(X_1)}
\]

\vspace{0.2cm}
\textbf{Dirección del sesgo:}
\begin{itemize}
  \item Si $\beta_2 > 0$ y $\text{Cov}(X_1, X_2) > 0$ $\Rightarrow$ sesgo positivo (sobre-estimación)
  \item Si $\beta_2 > 0$ y $\text{Cov}(X_1, X_2) < 0$ $\Rightarrow$ sesgo negativo (sub-estimación)
\end{itemize}
\end{frame}

\begin{frame}[fragile]{Ejemplo emp\'irico: efecto de IES privada}
\textbf{Modelo simple (omite ESTRATO):}
\begin{lstlisting}
modelo_omitido <- lm(MOD_INGLES_PUNT ~ PRIVADA, data = icfes_ni)
coef(modelo_omitido)

# (Intercept)     PRIVADA
#      154.82        7.35
\end{lstlisting}

\vspace{0.2cm}
\textbf{Modelo controlando por ESTRATO:}
\begin{lstlisting}
modelo_con_estrato <- lm(MOD_INGLES_PUNT ~ PRIVADA + ESTRATO,
                         data = icfes_ni)
coef(modelo_con_estrato)

# (Intercept)     PRIVADA      ESTRATO
#      148.34        4.12         2.87
\end{lstlisting}
\end{frame}

\begin{frame}{Ejemplo emp\'irico: efecto de IES privada (cont.)}
\textbf{Observaci\'on:} El coeficiente de PRIVADA cae de 7.35 a 4.12 al controlar por ESTRATO. El efecto aparente de IES privada estaba \textbf{inflado} por el sesgo de variable omitida.

\vspace{0.3cm}
\begin{alertblock}{Direcci\'on del sesgo}
\begin{itemize}
  \item ESTRATO correlacionado positivamente con PRIVADA ($\text{Cov} > 0$)
  \item ESTRATO tiene efecto positivo en ingl\'es ($\beta_2 > 0$)
  \item $\Rightarrow$ Sesgo positivo: $\hat{\gamma}_1 > \beta_1$ (sobre-estimaci\'on)
\end{itemize}
\end{alertblock}
\end{frame}

\begin{frame}[fragile]{Verificar la correlaci\'on entre PRIVADA y ESTRATO}
\begin{lstlisting}[basicstyle=\ttfamily\scriptsize]
# Correlacion entre PRIVADA y ESTRATO
cor(icfes_ni$PRIVADA, icfes_ni$ESTRATO)

# [1] 0.38  (correlacion positiva moderada)

# Tabla cruzada
table(icfes_ni$PRIVADA, icfes_ni$ESTRATO)

#        1   2   3   4   5   6
#   0  456 823 512 387 123  34  (publicas)
#   1  234 512 687 823 412  82  (privadas)

# Media de estrato por tipo IES
tapply(icfes_ni$ESTRATO, icfes_ni$PRIVADA, mean)

#        0        1
#     2.43     3.21

# Los estudiantes de IES privadas tienen estrato mas alto
\end{lstlisting}
\vspace{-2pt}
{\footnotesize \textbf{Conclusi\'on:} ESTRATO est\'a correlacionado con PRIVADA, y afecta el puntaje de ingl\'es. Omitir ESTRATO genera OVB.}
\end{frame}

\begin{frame}{Implicaciones del OVB}
\begin{alertblock}{Problema para inferencia causal}
\small Si queremos estimar el \textbf{efecto causal} de IES privada sobre ingl\'es, debemos controlar por \textbf{todas} las variables confusoras. De lo contrario, el coeficiente NO representa el efecto causal.
\end{alertblock}

\vspace{0.1cm}
\textbf{Estrategias para mitigar OVB:}
\begin{itemize}
  \item Incluir todas las variables relevantes (si est\'an disponibles)
  \item Usar variables instrumentales (IV)
  \item Dise\~nos experimentales o cuasi-experimentales (RCT, DiD, RD)
\end{itemize}

\vspace{0.1cm}
\begin{exampleblock}{Regla de oro}
\small Nunca concluir causalidad solo con regresi\'on observacional. Si ven que ``usar iPad'' se asocia con mejores notas, preg\'untense: \textquestiondown qu\'e variable omitida podr\'ia explicar ambas cosas?
\end{exampleblock}
\end{frame}

% ============================================================================
\section{VIF y multicolinealidad}
% ============================================================================

\begin{frame}{\textquestiondown Qu\'e es la multicolinealidad?}
\begin{block}{Definici\'on}
La \textbf{multicolinealidad} ocurre cuando dos o m\'as predictores est\'an \textbf{altamente correlacionados} entre s\'i.
\end{block}

\vspace{0.1cm}
\textbf{Tipos:}
\begin{itemize}
  \item \textbf{Perfecta:} Una variable es combinaci\'on lineal exacta de otras (raro; R lo detecta autom\'aticamente)
  \item \textbf{Alta:} Correlaci\'on fuerte pero no perfecta (com\'un en datos observacionales)
\end{itemize}

\vspace{0.1cm}
\textbf{Consecuencias:}
\begin{itemize}
  \item Errores est\'andar \textbf{inflados}, intervalos de confianza muy amplios
  \item Coeficientes inestables, dificultad para detectar significancia individual
\end{itemize}
\end{frame}

\begin{frame}{\textquestiondown Qu\'e es la multicolinealidad? (cont.)}
\begin{exampleblock}{Analog\'ia}
Horas en TikTok y horas en Instagram miden casi lo mismo (screen time en redes). Si meten ambas en el modelo, hay multicolinealidad: el modelo no puede separar el efecto de cada una.
\end{exampleblock}

\vspace{0.3cm}
\begin{alertblock}{Importante}
La multicolinealidad \textbf{no afecta} $R^2$ ni la capacidad de predicci\'on del modelo. Solo afecta la interpretaci\'on de los coeficientes individuales.
\end{alertblock}
\end{frame}

\begin{frame}{Factor de Inflación de la Varianza (VIF)}
\begin{block}{Definición}
El \textbf{VIF} (Variance Inflation Factor) mide cuánto se infla la varianza de $\hat{\beta}_j$ debido a la correlación con otros predictores.

Para cada predictor $X_j$:
\[
\text{VIF}_j = \frac{1}{1 - R_j^2}
\]
donde $R_j^2$ es el $R^2$ de la regresión auxiliar:
\[
X_j = \alpha_0 + \alpha_1 X_1 + \cdots + \alpha_{j-1} X_{j-1} + \alpha_{j+1} X_{j+1} + \cdots + \alpha_p X_p + e
\]

Es decir, regresamos $X_j$ sobre todos los demás predictores.
\end{block}

\vspace{0.2cm}
\textbf{Interpretación:}
\begin{itemize}
  \item VIF = 1 $\Rightarrow$ $X_j$ no está correlacionado con otros predictores
  \item VIF = 5 $\Rightarrow$ la varianza de $\hat{\beta}_j$ es 5 veces mayor que si $X_j$ fuera independiente
\end{itemize}
\end{frame}

\begin{frame}{Reglas de pulgar para VIF}
\textbf{Umbrales comunes:}
\begin{itemize}
  \item VIF $< 5$ $\Rightarrow$ multicolinealidad \textbf{baja} (aceptable)
  \item $5 \leq$ VIF $< 10$ $\Rightarrow$ multicolinealidad \textbf{moderada} (preocupante)
  \item VIF $\geq 10$ $\Rightarrow$ multicolinealidad \textbf{severa} (requiere acción)
\end{itemize}

\vspace{0.3cm}
\begin{alertblock}{Soluciones si VIF es alto}
\begin{enumerate}
  \item \textbf{Eliminar una de las variables correlacionadas} (la menos importante teóricamente)
  \item \textbf{Combinar variables:} crear un índice o promedio
  \item \textbf{Componentes principales (PCA):} reducir dimensionalidad
  \item \textbf{Regularización:} Ridge regression (penaliza coeficientes grandes)
  \item \textbf{Aumentar el tamaño muestral} (reduce errores estándar en general)
\end{enumerate}
\end{alertblock}
\end{frame}

\begin{frame}[fragile]{Calcular VIF en R}
\begin{lstlisting}
# Instalar paquete car (si no esta instalado)
# install.packages("car")
library(car)

# Calcular VIF para el modelo completo
vif(modelo_completo)

#  INGLES_SABER11_PUNT              ESTRATO
#                 1.23                 2.87
#              PRIVADA         LOG_INGRESOS
#                 3.45                 4.12
#        REGION_CARIBE      REGION_PACIFICA
#                 1.18                 1.09
#     REGION_ORINOQUIA        REGION_BOGOTA
#                 1.05                 1.32

# Todos los VIF < 5 => multicolinealidad baja
\end{lstlisting}

\vspace{0.1cm}
{\small \textbf{Conclusi\'on:} No hay problemas serios de multicolinealidad en este modelo.}
\end{frame}

\begin{frame}[fragile]{Ejemplo con multicolinealidad severa}
\begin{lstlisting}[basicstyle=\ttfamily\scriptsize]
# Crear una variable casi perfectamente correlacionada
icfes_ni$ESTRATO_x2 <- icfes_ni$ESTRATO * 2 + rnorm(nrow(icfes_ni), 0, 0.1)

cor(icfes_ni$ESTRATO, icfes_ni$ESTRATO_x2)

# [1] 0.9998  (correlacion casi perfecta)

# Ajustar modelo con ambas
modelo_colin <- lm(MOD_INGLES_PUNT ~ ESTRATO + ESTRATO_x2,
                   data = icfes_ni)

vif(modelo_colin)

#       ESTRATO    ESTRATO_x2
#       5432.12       5432.12

# VIF extremadamente alto => multicolinealidad severa
# Los errores estandar seran enormes
\end{lstlisting}

\vspace{0.1cm}
{\small \textbf{Soluci\'on:} Eliminar una de las dos variables (son redundantes).}
\end{frame}

% ============================================================================
\section{Selección de modelos: AIC y BIC}
% ============================================================================

\begin{frame}{Criterios de informaci\'on: AIC y BIC}
\begin{block}{Problema}
Queremos elegir el \textbf{mejor} modelo entre varios candidatos. \textquestiondown C\'omo balanceamos ajuste y complejidad?
\end{block}

\vspace{0.1cm}
\textbf{Dos criterios populares:}
\begin{enumerate}
  \item \textbf{AIC} (Akaike Information Criterion):
  \[
  \text{AIC} = n \ln\!\left(\frac{\text{SSE}}{n}\right) + 2(p + 1)
  \]

  \item \textbf{BIC} (Bayesian Information Criterion):
  \[
  \text{BIC} = n \ln\!\left(\frac{\text{SSE}}{n}\right) + (p + 1) \ln(n)
  \]
\end{enumerate}
\vspace{-2pt}
{\small \textbf{Interpretaci\'on:} Menor AIC/BIC $\Rightarrow$ mejor modelo. BIC penaliza m\'as la complejidad que AIC.}
\end{frame}

\begin{frame}{AIC vs BIC}
\textbf{AIC:}
\begin{itemize}
  \item Derivado de teor\'ia de informaci\'on (Kullback-Leibler divergence)
  \item Penaliza con $2(p+1)$; favorece modelos m\'as complejos
  \item Mejor para predicci\'on
\end{itemize}

\vspace{0.1cm}
\textbf{BIC:}
\begin{itemize}
  \item Derivado de teor\'ia bayesiana
  \item Penaliza con $(p+1) \ln(n)$; si $n > 8$, BIC penaliza m\'as que AIC
  \item Favorece modelos m\'as simples (parsimonia); mejor para inferencia
\end{itemize}

\vspace{0.2cm}
\begin{exampleblock}{Regla pr\'actica}
\begin{itemize}
  \item Si el objetivo es \textbf{predicci\'on}: usar AIC
  \item Si el objetivo es \textbf{inferencia/interpretaci\'on}: usar BIC
\end{itemize}
\end{exampleblock}
\end{frame}

\begin{frame}[fragile]{Calcular AIC y BIC en R}
\begin{lstlisting}[basicstyle=\ttfamily\scriptsize]
# Modelo 1: solo Saber 11
modelo1 <- lm(MOD_INGLES_PUNT ~ INGLES_SABER11_PUNT,
              data = icfes_ni)

# Modelo 2: + estrato y privada
modelo2 <- lm(MOD_INGLES_PUNT ~ INGLES_SABER11_PUNT +
                ESTRATO + PRIVADA,
              data = icfes_ni)

# Modelo 3: modelo completo (8 variables)
# (ya ajustado: modelo_completo)

# Comparar AIC
AIC(modelo1, modelo2, modelo_completo)

#                   df      AIC
# modelo1            3  42345.2
# modelo2            5  41823.6
# modelo_completo   10  41234.8

# Menor AIC => modelo_completo es el mejor segun AIC
\end{lstlisting}
\end{frame}

\begin{frame}[fragile]{Comparar con BIC}
\begin{lstlisting}[basicstyle=\ttfamily\scriptsize]
# Comparar BIC
BIC(modelo1, modelo2, modelo_completo)

#                   df      BIC
# modelo1            3  42367.1
# modelo2            5  41856.8
# modelo_completo   10  41302.4

# Menor BIC => modelo_completo tambien es el mejor segun BIC

# Diferencia AIC
AIC(modelo2) - AIC(modelo_completo)

# [1] 588.8  (modelo_completo es 589 puntos mejor en AIC)

# Diferencia BIC
BIC(modelo2) - BIC(modelo_completo)

# [1] 554.4
\end{lstlisting}
\vspace{-2pt}
{\footnotesize \textbf{Conclusi\'on:} El modelo completo es sustancialmente mejor que los modelos m\'as simples seg\'un ambos criterios.}
\end{frame}

\begin{frame}[fragile]{Selecci\'on autom\'atica con \texttt{step()}}
\begin{lstlisting}[basicstyle=\ttfamily\scriptsize]
# Seleccion automatica hacia atras (backward)
# Inicia con el modelo completo y va eliminando variables
modelo_backward <- step(modelo_completo, direction = "backward")

# Muestra los pasos:
# Step 1: elimina REGION_PACIFICA (no significativa)
# Step 2: se detiene (todas las demas son significativas)

# Modelo final
summary(modelo_backward)

# Puede hacerse tambien forward (hacia adelante) o both (ambos)
# backward: inicia con todas las variables, elimina una a una
# forward: inicia sin variables, agrega una a una
# both: combinacion de ambos
\end{lstlisting}
\end{frame}

\begin{frame}{Selecci\'on autom\'atica con \texttt{step()} (cont.)}
\begin{alertblock}{Precauci\'on}
La selecci\'on autom\'atica puede llevar a overfitting y problemas de inferencia (p-valores inflados). \'Usela como exploraci\'on, pero siempre valide con teor\'ia y conocimiento del problema.
\end{alertblock}

\vspace{0.3cm}
\begin{block}{Tipos de selecci\'on}
\begin{itemize}
  \item \textbf{backward:} inicia con todas las variables, elimina una a una
  \item \textbf{forward:} inicia sin variables, agrega una a una
  \item \textbf{both:} combinaci\'on de ambos (stepwise)
\end{itemize}
\end{block}
\end{frame}

% ============================================================================
\section{Interacciones}
% ============================================================================

\begin{frame}{\textquestiondown Qu\'e es una interacci\'on?}
\begin{block}{Definici\'on}
Una \textbf{interacci\'on} ocurre cuando el efecto de una variable $X_1$ sobre $Y$ \textbf{depende del nivel} de otra variable $X_2$.
\end{block}

\vspace{0.2cm}
\textbf{Modelo con interacci\'on:}
\[
Y = \beta_0 + \beta_1 X_1 + \beta_2 X_2 + \beta_3 (X_1 \times X_2) + \varepsilon
\]

\vspace{0.1cm}
\textbf{Interpretaci\'on:}
\begin{itemize}
  \item $\beta_1$ = efecto de $X_1$ cuando $X_2 = 0$
  \item $\beta_2$ = efecto de $X_2$ cuando $X_1 = 0$
  \item $\beta_3$ = c\'omo cambia el efecto de $X_1$ cuando $X_2$ aumenta en 1 unidad
\end{itemize}
\end{frame}

\begin{frame}{\textquestiondown Qu\'e es una interacci\'on? (cont.)}
\begin{exampleblock}{Ejemplo conceptual}
\textquestiondown El efecto de IES privada sobre ingl\'es es el \textbf{mismo} para todos los estratos, o es \textbf{mayor} en estratos altos? Similar a preguntarse: \textquestiondown el efecto de usar iPad en las notas es igual para todos, o depende de si adem\'as tienen acceso a internet en casa?
\end{exampleblock}
\end{frame}

\begin{frame}{Interpretación matemática de la interacción}
Modelo con interacción:
\[
Y = \beta_0 + \beta_1 X_1 + \beta_2 X_2 + \beta_3 (X_1 X_2) + \varepsilon
\]

Reordenando:
\[
Y = \beta_0 + (\beta_1 + \beta_3 X_2) X_1 + \beta_2 X_2 + \varepsilon
\]

\vspace{0.2cm}
El \textbf{efecto marginal de $X_1$} es:
\[
\frac{\partial Y}{\partial X_1} = \beta_1 + \beta_3 X_2
\]

\vspace{0.2cm}
\textbf{Interpretación:}
\begin{itemize}
  \item Si $\beta_3 = 0$ $\Rightarrow$ no hay interacción (efecto constante)
  \item Si $\beta_3 > 0$ $\Rightarrow$ el efecto de $X_1$ aumenta cuando $X_2$ aumenta
  \item Si $\beta_3 < 0$ $\Rightarrow$ el efecto de $X_1$ disminuye cuando $X_2$ aumenta
\end{itemize}
\end{frame}

\begin{frame}[fragile]{Ejemplo: interacción PRIVADA × ESTRATO}
\textbf{Pregunta:} ¿El efecto de IES privada varía según el estrato?

\begin{lstlisting}
# Modelo con interaccion
modelo_interaccion <- lm(MOD_INGLES_PUNT ~
                           PRIVADA * ESTRATO,
                         data = icfes_ni)

summary(modelo_interaccion)

# Coefficients:
#                  Estimate Std. Error t value Pr(>|t|)
# (Intercept)       149.23       1.12  133.24  < 2e-16 ***
# PRIVADA             2.34       1.87    1.25  0.21123
# ESTRATO             2.67       0.38    7.03  2.3e-12 ***
# PRIVADA:ESTRATO     0.89       0.52    1.71  0.08734 .
#
# Residual standard error: 19.45 on 5081 DF
# Multiple R-squared:  0.3512, Adjusted R-squared:  0.3508
\end{lstlisting}
\end{frame}

\begin{frame}{Interpretación del modelo con interacción}
Ecuación estimada:
\[
\widehat{\text{MOD\_INGLES}} = 149.23 + 2.34 \times \text{PRIVADA} + 2.67 \times \text{ESTRATO} + 0.89 \times (\text{PRIVADA} \times \text{ESTRATO})
\]

\vspace{0.2cm}
\textbf{Efecto de PRIVADA según ESTRATO:}
\[
\frac{\partial \text{Inglés}}{\partial \text{PRIVADA}} = 2.34 + 0.89 \times \text{ESTRATO}
\]

\begin{itemize}
  \item Estrato 1: efecto = $2.34 + 0.89 \times 1 = 3.23$ puntos
  \item Estrato 3: efecto = $2.34 + 0.89 \times 3 = 5.01$ puntos
  \item Estrato 5: efecto = $2.34 + 0.89 \times 5 = 6.79$ puntos
\end{itemize}

\vspace{0.2cm}
\textbf{Interpretación:} El efecto de IES privada \textbf{aumenta} con el estrato (0.89 puntos adicionales por cada nivel de estrato), aunque la interacción es solo marginalmente significativa (p = 0.087).
\end{frame}

\begin{frame}[fragile]{Visualizaci\'on de la interacci\'on}
\begin{lstlisting}[basicstyle=\ttfamily\scriptsize]
# Crear grafico de interaccion
library(ggplot2)

# Predecir para diferentes combinaciones
predicciones <- expand.grid(
  PRIVADA = c(0, 1),
  ESTRATO = 1:6
)
predicciones$MOD_INGLES_PUNT <- predict(modelo_interaccion,
                                         newdata = predicciones)

ggplot(predicciones, aes(x = ESTRATO, y = MOD_INGLES_PUNT,
                         color = factor(PRIVADA), group = PRIVADA)) +
  geom_line(size = 1.2) +
  geom_point(size = 3) +
  labs(title = "Interaccion PRIVADA x ESTRATO",
       x = "Estrato", y = "Puntaje ingles predicho",
       color = "Tipo IES") +
  scale_color_manual(values = c("blue", "red"),
                     labels = c("Publica", "Privada")) +
  theme_minimal()
\end{lstlisting}
\end{frame}

\begin{frame}[fragile]{Interacción entre dos variables continuas}
\begin{lstlisting}
# Ejemplo: interaccion INGLES_SABER11 x LOG_INGRESOS
modelo_cont_int <- lm(MOD_INGLES_PUNT ~
                        INGLES_SABER11_PUNT * LOG_INGRESOS,
                      data = icfes_ni)

summary(modelo_cont_int)

# Coefficients:
#                                  Estimate Std. Error t value Pr(>|t|)
# (Intercept)                       22.345      3.821    5.85  5.2e-09 ***
# INGLES_SABER11_PUNT                1.523      0.062   24.56  < 2e-16 ***
# LOG_INGRESOS                       1.234      0.512    2.41   0.0160 *
# INGLES_SABER11_PUNT:LOG_INGRESOS   0.023      0.009    2.56   0.0105 *
#
# La interaccion es significativa (p = 0.01)
\end{lstlisting}

\vspace{0.2cm}
\textbf{Interpretación:} El efecto de Saber 11 sobre Saber Pro es \textbf{mayor} en estudiantes de hogares con mayores ingresos.
\end{frame}

% ============================================================================
\section{Resumen}
% ============================================================================

\begin{frame}{Resumen de regresión múltiple}
\begin{block}{Conceptos clave}
\begin{itemize}
  \item Modelo: $Y = \beta_0 + \beta_1 X_1 + \cdots + \beta_p X_p + \varepsilon$
  \item $\beta_j$ = efecto de $X_j$ \textbf{controlando por} las demás variables
  \item OLS minimiza $\sum e_i^2$ para estimar $\hat{\beta}_0, \ldots, \hat{\beta}_p$
  \item $R^2_{\text{adj}}$ penaliza por número de variables
  \item AIC/BIC balancean ajuste y complejidad
\end{itemize}
\end{block}

\begin{block}{Pruebas de significancia}
\begin{itemize}
  \item Prueba F global: $H_0: \beta_1 = \cdots = \beta_p = 0$
  \item Prueba t individual: $H_0: \beta_j = 0$ (controlando por los demás)
  \item Reportar ambas + IC 95\%
\end{itemize}
\end{block}
\end{frame}

\begin{frame}{Problemas comunes y soluciones}
\begin{block}{Sesgo por variable omitida (OVB)}
\begin{itemize}
  \item Omitir variables relevantes correlacionadas con $X$ sesga $\hat{\beta}$
  \item Soluci\'on: incluir todas las variables confusoras
\end{itemize}
\end{block}

\begin{block}{Multicolinealidad}
\begin{itemize}
  \item Predictores altamente correlacionados $\Rightarrow$ errores est\'andar inflados
  \item Detectar con VIF ($> 5$ preocupante, $> 10$ severo)
  \item Soluci\'on: eliminar variables redundantes, PCA, regularizaci\'on
\end{itemize}
\end{block}

\begin{block}{Interacciones}
\begin{itemize}
  \item El efecto de $X_1$ depende de $X_2$: $Y \sim X_1 \times X_2$
  \item $\beta_3$ = c\'omo cambia el efecto de $X_1$ con $X_2$
\end{itemize}
\end{block}
\end{frame}

\begin{frame}{Checklist para análisis de regresión múltiple}
\begin{enumerate}
  \item \textbf{Exploración inicial:} scatter plots, correlaciones
  \item \textbf{Ajustar modelo:} \texttt{lm(Y \textasciitilde{} X1 + X2 + ...)}
  \item \textbf{Evaluar ajuste:} $R^2$, $R^2_{\text{adj}}$, prueba F global
  \item \textbf{Significancia individual:} pruebas t, p-valores, IC 95\%
  \item \textbf{Diagnóstico:}
    \begin{itemize}
      \item Residuos: \texttt{plot(modelo)}
      \item Multicolinealidad: \texttt{vif(modelo)}
      \item Normalidad: Shapiro-Wilk, Q-Q plot
      \item Homocedasticidad: Breusch-Pagan test
    \end{itemize}
  \item \textbf{Selección de modelo:} AIC/BIC, validación cruzada
  \item \textbf{Interacciones:} explorar si los efectos son condicionales
  \item \textbf{Reportar:} tabla de coeficientes, IC, $R^2$, interpretación
\end{enumerate}
\end{frame}

\begin{frame}{Ejemplo final: reporte profesional}
\textbf{Objetivo:} Modelar el puntaje de inglés en Saber Pro para estudiantes de Negocios Internacionales.

\vspace{0.2cm}
\textbf{Variables:}
\begin{itemize}
  \item Dependiente: MOD\_INGLES\_PUNT
  \item Independientes: Saber 11, estrato, tipo IES, log(ingresos), región
\end{itemize}

\vspace{0.2cm}
\textbf{Resultados:}
\begin{itemize}
  \item $R^2_{\text{adj}} = 0.714$ (71.4\% de variabilidad explicada)
  \item Prueba F: $F = 1586$, p $< 0.001$ (modelo altamente significativo)
  \item Todos los predictores significativos excepto Región Pacífica
  \item VIF $< 5$ para todas las variables (sin multicolinealidad severa)
  \item Residuos aproximadamente normales (Shapiro p = 0.08)
  \item Leve heterocedasticidad (Breusch-Pagan p = 0.03), considerar errores robustos
\end{itemize}

\vspace{0.2cm}
\textbf{Conclusión:} El modelo ajusta bien y revela efectos significativos de Saber 11, estrato, tipo de IES, ingresos y región sobre el desempeño en inglés.
\end{frame}

\begin{frame}{Pr\'oximos pasos}
\begin{exampleblock}{Temas avanzados (fuera del alcance de este curso)}
\begin{itemize}
  \item Diagn\'osticos avanzados: influential points (Cook's distance), leverage
  \item Errores est\'andar robustos (HC, HAC)
  \item Modelos no lineales: polinomios, splines, GAM
  \item Regresi\'on regularizada: Ridge, Lasso, Elastic Net
  \item Regresi\'on log\'istica (para variables dependientes binarias)
\end{itemize}
\end{exampleblock}

\begin{block}{Recursos recomendados}
\begin{itemize}
  \item Wooldridge, J. (2019). \textit{Introductory Econometrics}. Cap\'itulos 3-7.
  \item James, G. et al. (2021). \textit{An Introduction to Statistical Learning}. Cap\'itulo 3.
  \item R Documentation: \texttt{?lm}, \texttt{?summary.lm}, \texttt{?vif}
\end{itemize}
\end{block}
\end{frame}

\end{document}
